% \iffalse meta-comment
%
% Copyright (C) 2017- by Yanshuo Chu <yanshuoc@gmail.com>
%
% This file may be distributed and/or modified under the
% conditions of the LaTeX Project Public License, either version 1.3a
% of this license or (at your option) any later version.
% The latest version of this license is in:
%
% http://www.latex-project.org/lppl.txt
%
% and version 1.3a or later is part of all distributions of LaTeX
% version 2004/10/01 or later.
%
% \fi
%
% \section{终稿的术语与符号表}
%
% \changes{v3.2a}{2026/09/13}{中英两线统一到 \pkg{glossaries-extra}，后端
%   \cs{makenoidxglossaries}；缩略语表改成两列悬挂与符号表同款，加 \cs{abbr} 一族与
%   索引挂钩 \cs{hitindex}；中文索引的拼音排序键查 \pkg{xpinyin} 自动推（只在
%   \hologo{XeLaTeX} 下装），排序键用实际符；缩写的组尾边界跨版本对齐}
%
% ^^A ======================================================================
% ^^A Module final-glossary: 术语与符号表（glossaries 配置 + 示例条目）（hit-thesis）
% ^^A ======================================================================
%    \begin{macrocode}
%<*final-glossary>
%    \end{macrocode}
%
% \changes{v3.1d}{2025/03/03}{将 \cs{glossarystyle} 替换为 \cs{setglossarystyle}。\issue[200]}
% 在 \pkg{glossaries} 中表示 \cs{glossarystle} 这个命令 Deprecated in v3.08a. Removed in v4.50.
% 两条语言线原先各装各的包：英文装 \pkg{glossaries-extra} 用
% \cs{setabbreviationstyle}，中文装 \pkg{glossaries} 用 \cs{setacronymstyle}，
% 后者是上一代接口，\pkg{glossaries-extra} 里已经改成报错。要把缩略语表并成
% 一个部件，先得统一到 extra。
%
% 后端两侧都用 \cs{makenoidxglossaries}，纯 \hologo{TeX} 排序，不调外部程序。
% 中文侧原先调 \cs{makeglossaries} 却从没有任何东西 \cs{printglossary}，
% 每次编译写出一个没人读的 \file{.glo} 再吃一条警告，一并撤掉。
%    \begin{macrocode}
\RequirePackage [ xindy , order=letter ] { glossaries-extra }
%    \end{macrocode}
%
% \pkg{xpinyin} 只为中文索引的拼音排序键而装，见下文。它随包带一张四万七千行
% 的汉字拼音表，纯 \hologo{TeX} 查表，不调外部程序，加载实测四十毫秒。
%
% 只有 \hologo{XeLaTeX} 装得了它——\pkg{xpinyin} 见到 \hologo{LuaLaTeX} 直接
% 报 Critical Error 停机。查表宏 \cs{\_\_hit\_pinyin\_lookup:nN} 本来就以
% \cs{cs\_if\_exist:c} 探每个字的表项，包不在等于全部查不到，走的正是「查不到
% 就提示手写 \option{pinyin} 排序键」那条既有的路，不必另写分支。
%    \begin{macrocode}
\sys_if_engine_luatex:F
  { \RequirePackage { xpinyin } }
%    \end{macrocode}
%
% 样式代码存进 \pkg{glossaries} 时就已经记号化完了，写在里面的
% \cs{ExplSyntaxOn} 不起作用。环境正文本来就是排版内容，直接用 2e 写法。
%    \begin{macrocode}
%    \end{macrocode}
%
% 三列的宽度从 \cs{glsdescwidth} 分，不另写死数值：\pkg{glossaries} 按版心算出
% 这个长度，两列的 \texttt{engstyle} 用的就是它，照它切开总宽才不会超版心。
% 写死过一版 \texttt{2.5cm+6cm+\cs{glsdescwidth}}，中文变体每本都多一条
% Overfull（\cs{make} \texttt{logcheck} 当场拦下）。
%
% 首列前置 \cs{leavevmode}。首列的内容以 \cs{glstarget} 开头，它经 \pkg{hyperref}
% 的 \cs{Hy@raisedlink} 在竖直模式下先放一个零高盒，于是 \texttt{p} 列的
% \cs{vtop} 首项不是正文行而是那个盒（或它带出的胶），\cs{vtop} 的高按首项算，
% 行高就跟着变。\TeX~Live 2024 与 2025 的内核在这里处理不同，同一张表两版差
% 1.1\,pt，整行错开 3 个像素。先 \cs{leavevmode} 起段，锚点落进水平表，首项恒为
% 正文行，三版一致。
%    \begin{macrocode}
%    \end{macrocode}
%
% \emph{缩略语表与符号表同一套版式。}CTAN 上有这类页的中文模板五家里四家
% 是两列 description 列表（thu 的 denotation、bit 的 symbols、nju 与 fdu 的
% notation），说明一列排到右边距、折行悬挂；中英全称合成一格“体细胞拷贝数
% 变异（Somatic copy number alteration）”。先前那张三列长表（中文全称单占
% 一列）右边必然空一截，且与符号表两个长相，合并成一页时列也对不齐。改走
% \env{hit@symbols} 同一只列表（labelwidth 2.5\,cm、labelsep 0.5\,cm），
% 合并页两段的说明列于是天然对齐；\option{labelwidth} 键两段都认。
% 英文版说明只排英文全称（没有就中文），括号档位跟正文首次展开同一套
% （中文版全角）。
%    \begin{macrocode}
\cs_new:Npn \__hit_abbr_desc:n #1
  {
    \tl_if_blank:eTF { \hitabbrzh {#1} }
      { \hitabbren {#1} }
      {
        \tl_if_blank:eTF { \hitabbren {#1} }
          { \hitabbrzh {#1} }
          { \hitabbrzh {#1} （ \hitabbren {#1} ） }
      }
  }
\cs_new:Npn \__hit_abbr_desc_en:n #1
  {
    \tl_if_blank:eTF { \hitabbren {#1} }
      { \hitabbrzh {#1} }
      { \hitabbren {#1} }
  }
\newglossarystyle{hitabbrstyle}{%
  \RenewDocumentEnvironment{theglossary}{}%
    {\begin{hit@symbols}[labelwidth=\l__hit_symbols_width_tl]}%
    {\end{hit@symbols}}%
  \renewcommand*{\glossaryheader}{}%
  \renewcommand*{\glsgroupheading}[1]{}%
  \renewcommand*{\glsgroupskip}{}%
  \renewcommand*{\glossentry}[2]{%
    \item[\glsentryitem{##1}\glstarget{##1}{\glossentryname{##1}}]%
    \__hit_abbr_desc:n {##1}%
  }%
  \renewcommand*{\subglossentry}[3]{\glossentry{##2}{##3}}%
}
\newglossarystyle{engstyle}{%
  \RenewDocumentEnvironment{theglossary}{}%
    {\begin{hit@symbols}[labelwidth=\l__hit_symbols_width_tl]}%
    {\end{hit@symbols}}%
  \renewcommand*{\glossaryheader}{}%
  \renewcommand*{\glsgroupheading}[1]{}%
  \renewcommand*{\glsgroupskip}{}%
  \renewcommand*{\glossentry}[2]{%
    \item[\glsentryitem{##1}\glstarget{##1}{\glossentryname{##1}}]%
    \__hit_abbr_desc_en:n {##1}%
  }%
  \renewcommand*{\subglossentry}[3]{\glossentry{##2}{##3}}%
}
%    \end{macrocode}
%
% \cs{newacronym} 是 v3.1e 的写法，渲染冻结在那一版：中文侧首次只排全称、
% 之后排短形（v3.1e 靠 \cs{setacronymstyle}\marg{short-long} 再把
% \cs{genacrfullformat} 改成只取 \cs{glsentrylong} 凑出来，上游现成的
% \texttt{long-only-short-only} 就是这个行为）；英文侧照旧 \texttt{long-short}。
% \cs{genacrfullformat} 那个改写在 extra 下已经不起作用，留着是误导，删掉。
%    \begin{macrocode}
\bool_if:NTF \g__hit_en_bool
  { \setabbreviationstyle [ acronym ] { long-short } }
  { \setabbreviationstyle [ acronym ] { long-only-short-only } }
%    \end{macrocode}
%
% \pkg{glossaries-extra} 把 \cs{setacronymstyle} 做成了硬报错，指向
% \cs{setabbreviationstyle}。v3.1e 的用户可能在自己的 \file{.sty} 里写过它，
% 直接报错会让老文档编不过，所以拦下来转发，只提示一次。
%    \begin{macrocode}
\RenewDocumentCommand \setacronymstyle { o m }
  {
    \hit@warn@renamed@command { setacronymstyle }
      { setabbreviationstyle[acronym] }
    \IfNoValueTF {#1}
      { \setabbreviationstyle [ acronym ] {#2} }
      { \setabbreviationstyle [ #1 ] {#2} }
  }
\makenoidxglossaries
%    \end{macrocode}
%
% \subsection{缩略语的声明与使用}
%
% \option{abbr} 族的每一条经这里落成一个 \pkg{glossaries-extra} 词条。
% 规范要求「文中第一次出现的缩写词应该用括号注明英文原词」，中文全称在前时
% 括号里依次是英文原词与缩写：
% \begin{center}
% 体细胞拷贝数变异（Somatic copy number alteration，SCNA）
% \end{center}
% 三种取值各一种排法，中英文版又各一套括号：
% \begin{longtblr}[caption={缩写首次出现的排法}, label={tab:impl-abbr}]
%   {colspec={l l l}}
% \toprule
% 定义 & 中文版首次 & 英文版首次 \\ \midrule
% \texttt{zh}+\texttt{en} & X（Y，S） & Y (S) \\
% 只有 \texttt{zh}       & X（S）    & X (S) \\
% 只有 \texttt{en}       & S（Y）    & Y (S) \\
% \bottomrule
% \end{longtblr}
%
% 落实的办法是在声明时就把整串首次形式算好，存进词条的 \texttt{long} 字段，
% 再配上游现成的 \texttt{long-only-short-only} 样式——它首次排 \texttt{long}、
% 之后排 \texttt{short}，正是要的行为。复数另存 \texttt{longplural}，
% 括号里的缩写取复数而全称不动，中文全称加不上 \texttt{s}。
%
% v3.1e 是把括号手写进 \cs{newacronym} 的第三个参数里，单数看着对而复数漏了，
% 排出来半角括号套全角括号。现在括号由样式生成，两边不会再各走各的。
%    \begin{macrocode}
\tl_new:N \l__hit_abbr_zh_tl
\tl_new:N \l__hit_abbr_en_tl
\tl_new:N \l__hit_abbr_pinyin_tl
\tl_new:N \l__hit_abbr_label_tl
\tl_new:N \l__hit_abbr_long_tl
\tl_new:N \l__hit_abbr_longpl_tl
\tl_new:N \l__hit_abbr_pick_tl
\keys_define:nn { hit / abbr / entry }
  {
    zh     .tl_set:N = \l__hit_abbr_zh_tl ,
    en     .tl_set:N = \l__hit_abbr_en_tl ,
    pinyin .tl_set:N = \l__hit_abbr_pinyin_tl ,
    label  .tl_set:N = \l__hit_abbr_label_tl ,
  }
%    \end{macrocode}
%
% 判断一串记号里有没有汉字。逐记号取字符码，超过 \texttt{U+2E7F} 就算——
% 那一段之前是拉丁、希腊、标点，之后是部首与汉字。控制序列跳过，取不到字符码。
% 实测二十万次两秒半，单次十几微秒，而且只在声明时做一次，不进正文热路径。
%    \begin{macrocode}
\bool_new:N \l__hit_abbr_cjk_bool
\cs_new_protected:Npn \__hit_abbr_cjk_scan:n #1
  {
    \tl_map_inline:nn {#1}
      {
        \token_if_cs:NF ##1
          {
            \int_compare:nNnT { `##1 } > { "2E7F }
              { \bool_set_true:N \l__hit_abbr_cjk_bool }
          }
      }
  }
\cs_new_protected:Npn \__hit_abbr_cjk_set:n #1
  {
    \bool_set_false:N \l__hit_abbr_cjk_bool
    \__hit_abbr_cjk_scan:n {#1}
  }
\cs_generate_variant:Nn \__hit_abbr_cjk_set:n { V }
%    \end{macrocode}
%
% \begin{macro}{\hit@abbr@declare:nn}
% \meta{短形}\marg{取值}。取值不含等号就当裸值，按有没有汉字归 \option{zh}
% 或 \option{en}。
%    \begin{macrocode}
\cs_new_protected:Npn \hit@abbr@declare:nn #1#2
  {
    \tl_clear:N \l__hit_abbr_zh_tl
    \tl_clear:N \l__hit_abbr_en_tl
    \tl_clear:N \l__hit_abbr_pinyin_tl
    \tl_set:Nn \l__hit_abbr_label_tl {#1}
    \tl_if_in:nnTF {#2} { = }
      { \hit@keys@set:nn { hit / abbr / entry } {#2} }
      {
        \__hit_abbr_cjk_set:n {#2}
        \bool_if:NTF \l__hit_abbr_cjk_bool
          { \tl_set:Nn \l__hit_abbr_zh_tl {#2} }
          { \tl_set:Nn \l__hit_abbr_en_tl {#2} }
      }
    \__hit_abbr_compose:n {#1}
    \use:x
      {
        \exp_not:N \newabbreviation
          [
            longplural = { \exp_not:V \l__hit_abbr_longpl_tl } ,
            hitzh = { \exp_not:V \l__hit_abbr_zh_tl } ,
            hiten = { \exp_not:V \l__hit_abbr_en_tl } ,
            description = { \exp_not:V \l__hit_abbr_long_tl } ,
          ]
          { \exp_not:V \l__hit_abbr_label_tl }
          { \exp_not:n {#1} }
          { \exp_not:V \l__hit_abbr_long_tl }
      }
  }
%    \end{macrocode}
% \end{macro}
%
% 拼首次形式。中文版用全角括号与全角逗号，英文版用半角括号与空格。
%    \begin{macrocode}
%    \end{macrocode}
%
% 括号与逗号取哪一档由 \option{abbr}/\option{punct} 定，\texttt{auto} 按语种走。
% 排法本身仍按语种：中文版全称在前、英文版全称在前但括号里只有缩写，这个不受
% 标点档位影响，档位只换那三个符号。
%    \begin{macrocode}
\tl_new:N \l__hit_abbr_lparen_tl
\tl_new:N \l__hit_abbr_rparen_tl
\tl_new:N \l__hit_abbr_comma_tl
\cs_new_protected:Npn \__hit_abbr_punct_pick:
  {
    \bool_lazy_or:nnTF
      { \str_if_eq_p:Vn \g__hit_abbr_punct_str { half } }
      {
        \bool_lazy_and_p:nn
          { \str_if_eq_p:Vn \g__hit_abbr_punct_str { auto } }
          { \bool_if_p:N \g__hit_en_bool }
      }
      {
        \tl_set:Nn \l__hit_abbr_lparen_tl { \c_space_tl ( }
        \tl_set:Nn \l__hit_abbr_rparen_tl { ) }
        \tl_set:Nn \l__hit_abbr_comma_tl { ,~ }
      }
      {
        \tl_set:Nn \l__hit_abbr_lparen_tl { （ }
        \tl_set:Nn \l__hit_abbr_rparen_tl { ） }
        \tl_set:Nn \l__hit_abbr_comma_tl { ， }
      }
  }
\cs_new_protected:Npn \__hit_abbr_compose:n #1
  {
    \__hit_abbr_punct_pick:
    \bool_if:NTF \g__hit_en_bool
      { \__hit_abbr_compose_en:n {#1} }
      { \__hit_abbr_compose_zh:n {#1} }
  }
\cs_new_protected:Npn \__hit_abbr_compose_en:n #1
  {
    \tl_if_empty:NTF \l__hit_abbr_en_tl
      { \tl_set_eq:NN \l__hit_abbr_pick_tl \l__hit_abbr_zh_tl }
      { \tl_set_eq:NN \l__hit_abbr_pick_tl \l__hit_abbr_en_tl }
    \tl_set:Nx \l__hit_abbr_long_tl
      {
        \exp_not:V \l__hit_abbr_pick_tl
        \exp_not:V \l__hit_abbr_lparen_tl #1 \exp_not:V \l__hit_abbr_rparen_tl
      }
%    \end{macrocode}
%
% 全称的复数只在它是西文时才加 \texttt{s}。英文版里引一条中文全称的条目
% 是少见但合法的写法，那时候「树结构折筷过程s」显然不对。
%    \begin{macrocode}
    \__hit_abbr_cjk_set:V \l__hit_abbr_pick_tl
    \tl_set:Nx \l__hit_abbr_longpl_tl
      {
        \exp_not:V \l__hit_abbr_pick_tl
        \bool_if:NF \l__hit_abbr_cjk_bool { s }
        \exp_not:V \l__hit_abbr_lparen_tl #1 s \exp_not:V \l__hit_abbr_rparen_tl
      }
  }
\cs_new_protected:Npn \__hit_abbr_compose_zh:n #1
  {
    \tl_if_empty:NTF \l__hit_abbr_zh_tl
      {
        \tl_set:Nx \l__hit_abbr_long_tl
          {
            #1 \exp_not:V \l__hit_abbr_lparen_tl
            \exp_not:V \l__hit_abbr_en_tl \exp_not:V \l__hit_abbr_rparen_tl
          }
        \tl_set:Nx \l__hit_abbr_longpl_tl
          {
            #1 s \exp_not:V \l__hit_abbr_lparen_tl
            \exp_not:V \l__hit_abbr_en_tl \exp_not:V \l__hit_abbr_rparen_tl
          }
      }
      {
        \tl_set:Nx \l__hit_abbr_long_tl
          {
            \exp_not:V \l__hit_abbr_zh_tl \exp_not:V \l__hit_abbr_lparen_tl
            \tl_if_empty:NF \l__hit_abbr_en_tl
              {
                \exp_not:V \l__hit_abbr_en_tl
                \exp_not:V \l__hit_abbr_comma_tl
              }
            #1 \exp_not:V \l__hit_abbr_rparen_tl
          }
        \tl_set:Nx \l__hit_abbr_longpl_tl
          {
            \exp_not:V \l__hit_abbr_zh_tl \exp_not:V \l__hit_abbr_lparen_tl
            \tl_if_empty:NF \l__hit_abbr_en_tl
              {
                \exp_not:V \l__hit_abbr_en_tl
                \exp_not:V \l__hit_abbr_comma_tl
              }
            #1 s \exp_not:V \l__hit_abbr_rparen_tl
          }
      }
  }
\glsaddstoragekey { hitzh } { } { \hitabbrzh }
\glsaddstoragekey { hiten } { } { \hitabbren }
\setabbreviationstyle [ abbreviation ] { long-only-short-only }
%    \end{macrocode}
%
% \subsection{拼音排序键}
%
% 中文索引靠一个人工排序键分组，\pkg{splitidx} 底下是 \pkg{makeindex}，
% 它不认汉字。v3.1e 的写法是每处手写 \texttt{ti!体细胞拷贝数变异}。
%
% \pkg{xpinyin} 随包带一张四万七千行的表，每个汉字一条
% \cs{XPYU}\marg{字}\marg{码点}\marg{带调拼音}，存成
% \cs{c\_\_xpinyin\_\meta{码点}\_tl}。查表是纯 \hologo{TeX}，不调外部程序，
% 整张表加载实测四十毫秒。所以首字的拼音可以推出来，声调转成数字排在后面
% （\texttt{tǐ} 记作 \texttt{ti3}），同音节按调号排，与汉语词典一致。
%
% 两处必须提示用户：库里查不到，以及首字是多音字。两种都退回取库里的首选读音
% 并说明怎么手写覆盖。
%    \begin{macrocode}
\tl_new:N \l__hit_pinyin_tl
\tl_new:N \l__hit_pinyin_head_tl
\prop_new:N \c__hit_pinyin_tone_prop
\cs_new_protected:Npn \__hit_pinyin_tone_put:nn #1#2
  { \prop_put:Nnn \c__hit_pinyin_tone_prop {#1} {#2} }
\clist_map_inline:nn
  {
    {ā}{a1},{á}{a2},{ǎ}{a3},{à}{a4},
    {ē}{e1},{é}{e2},{ě}{e3},{è}{e4},
    {ī}{i1},{í}{i2},{ǐ}{i3},{ì}{i4},
    {ō}{o1},{ó}{o2},{ǒ}{o3},{ò}{o4},
    {ū}{u1},{ú}{u2},{ǔ}{u3},{ù}{u4},
    {ǖ}{v1},{ǘ}{v2},{ǚ}{v3},{ǜ}{v4},{ü}{v},
  }
  { \__hit_pinyin_tone_put:nn #1 }
%    \end{macrocode}
%
% 把带调拼音拆成「无调音节 + 调号」。逐字符查上表，查得到就换成基字母并记下
% 调号，查不到原样留下。轻声没有调号，补 \texttt{5}。
%    \begin{macrocode}
\tl_new:N \l__hit_pinyin_tone_tl
\cs_new_protected:Npn \__hit_pinyin_detone:n #1
  {
    \tl_clear:N \l__hit_pinyin_head_tl
    \tl_set:Nn \l__hit_pinyin_tone_tl { 5 }
    \tl_map_inline:nn {#1}
      {
        \prop_get:NnNTF \c__hit_pinyin_tone_prop {##1} \l_tmpa_tl
          {
            \tl_put_right:Nx \l__hit_pinyin_head_tl
              { \tl_head:V \l_tmpa_tl }
            \tl_if_empty:NF \l_tmpa_tl
              {
                \tl_set:Nx \l_tmpb_tl { \tl_tail:V \l_tmpa_tl }
                \tl_if_empty:NF \l_tmpb_tl
                  { \tl_set_eq:NN \l__hit_pinyin_tone_tl \l_tmpb_tl }
              }
          }
          { \tl_put_right:Nn \l__hit_pinyin_head_tl {##1} }
      }
    \tl_put_right:NV \l__hit_pinyin_head_tl \l__hit_pinyin_tone_tl
  }
%    \end{macrocode}
%
% \begin{macro}{\hit@pinyin@head:nN}
% \meta{词}\meta{存放处}：取首字的排序键。
%    \begin{macrocode}
\cs_new_protected:Npn \hit@pinyin@head:nN #1#2
  {
    \tl_set:Nx \l__hit_pinyin_tl { \tl_head:n {#1} }
    \exp_args:Nf \__hit_pinyin_lookup:nN
      { \tl_use:N \l__hit_pinyin_tl } #2
  }
\cs_new_protected:Npn \__hit_pinyin_lookup:nN #1#2
  {
    \cs_if_exist:cTF { c__xpinyin_ \int_value:w `#1 _tl }
      {
        \exp_args:Nv \__hit_pinyin_detone:n
          { c__xpinyin_ \int_value:w `#1 _tl }
        \tl_set_eq:NN #2 \l__hit_pinyin_head_tl
%    \end{macrocode}
%
% 提示里给的是排序键该写成的样子，也就是转过数字的形式，不是库里那个带调的。
%    \begin{macrocode}
        \clist_if_exist:cT { c__xpinyin_multiple_ \int_value:w `#1 _clist }
          {
            \msg_warning:nnxxx { hithesis } { pinyin-multiple } {#1}
              { \tl_use:N \l__hit_pinyin_head_tl }
              { \use:c { c__xpinyin_multiple_ \int_value:w `#1 _clist } }
          }
      }
      {
        \msg_warning:nnn { hithesis } { pinyin-missing } {#1}
        \tl_clear:N #2
      }
  }
%    \end{macrocode}
% \end{macro}
%
% \begin{macro}{\abbr,\abbrs}
% 正文里用缩写。\cs{abbr} 首次排全称、之后排缩写；\cs{abbrs} 是复数；
% 带星号是强制全写，不改「用过」状态，摘要或新一章想再展开一次时用。
%
% 缩写排出来的东西裹在 \pkg{glossaries} 的组里，组尾边界两版 xeCJK
% 各走各的：首次形态收在闭括号上，按“闭号接汉字摘胶”该贴排，2026 的
% 边界体系却隔着组把胶发了出来——尾后垫一只零宽盒把类对断掉，两版都
% 归到贴排；之后的形态收在西文字母上，接汉字要一道中西文胶，2026 隔组
% 发得出来、老版发不出来——置上“刚出组”标记，交给 hit-format 的
% 消费点补。分支跟 \pkg{glossaries-extra} 排版用的同一根轴走：排完
% 问 \cs{glsxtrifwasfirstuse}，刚才排的是首次（收在闭括号上）就垫盒
% 贴排，是缩写（收在字母上）就置标记补胶。全角标点哨兵在这儿靠不
% 住——链接收尾的 whatsit 会把它盖掉。
%
% 这一垫合的是 Word 的账：书写范例版式里“……，TSSBP）是”的括号到
% 汉字步进恰好一个字格（12.45\,bp，Word 实排实测），闭括号剪切贴排；
% 隔着链接组界两版 xeCJK 都发不出剪切，只有断掉类对才回到一格。
%
% 模板不暴露 \cs{gls}：概念表里没有 glossary 这一项，只有缩写、符号与索引。
%    \begin{macrocode}
\cs_new_protected:Npn \__hit_abbr_edge:nn #1#2
  {
    #2
    \glsxtrifwasfirstuse
      {
        \hbox:n { }
        \bool_if_exist:NT \g_docxlike_format_verb_exit_bool
          { \bool_gset_false:N \g_docxlike_format_verb_exit_bool }
      }
      { \cs_if_exist_use:N \docxlike_format_group_exit: }
  }
\NewDocumentCommand \abbr { s O{} m }
  {
    \__hit_abbr_use:nn {#2} {#3}
    \IfBooleanTF {#1} { \glsentrylong {#3} }
      { \__hit_abbr_edge:nn {#3} { \gls {#3} } }
  }
\NewDocumentCommand \abbrs { s O{} m }
  {
    \__hit_abbr_use:nn {#2} {#3}
    \IfBooleanTF {#1} { \glsentrylongpl {#3} }
      { \__hit_abbr_edge:nn {#3} { \glspl {#3} } }
  }
%    \end{macrocode}
% \end{macro}
%
% \begin{macro}{\hitindex}
% 索引与缩略语是两件事：索引回答「这个词在哪几页」，缩略语回答「这个缩写是什么」。
% 但要索引的词与要缩写的词重合度很高，所以缩写条目顺带登记索引，开关是
% \option{abbr}/\option{indexing}，单处可用可选参数压过。
%
% \cs{hitindex} 给不缩写的词：\cs{hitindex}\marg{zh=乔峰, en=Qiao~Feng} 排出当前
% 语种的那一份，两个索引都登记。裸值按有没有汉字归边。带星号只登记不排——正文里
% 写的是「乔帮主」而索引条目要「乔峰」时用。
%
% 中文索引的排序键取首字的拼音，见上面的 \cs{hit@pinyin@head:nN}；写了
% \option{pinyin} 就以写的为准。
%    \begin{macrocode}
\tl_new:N \l__hit_index_zh_tl
\tl_new:N \l__hit_index_en_tl
\tl_new:N \l__hit_index_pinyin_tl
\bool_new:N \l__hit_index_on_bool
\keys_define:nn { hit / index / one }
  {
    zh     .tl_set:N = \l__hit_index_zh_tl ,
    en     .tl_set:N = \l__hit_index_en_tl ,
    pinyin .tl_set:N = \l__hit_index_pinyin_tl ,
  }
\keys_define:nn { hit / abbr / use }
  {
    indexing .bool_set:N = \l__hit_index_on_bool ,
    indexing .default:n = true ,
  }
%    \end{macrocode}
%
% 真正往两个索引里写。中文那份要排序键，缺了就不写——没有键的条目在
% \pkg{makeindex} 里会排到莫名其妙的地方，不如不进。
%
% 键与词之间用的是 \texttt{@}（\pkg{makeindex} 的实际符：@ 前排序、@ 后
% 排印），从前误写成层级符 \texttt{!}，拼音键被当成父条目印了出来
% （“duan4”一行、词缩进在下），官方示例（规范附录 6）里只有词，
% 组标题是单个大写字母。
%    \begin{macrocode}
\cs_new_protected:Npn \__hit_index_emit:
  {
    \tl_if_empty:NF \l__hit_index_zh_tl
      {
        \tl_if_empty:NT \l__hit_index_pinyin_tl
          {
            \exp_args:NV \hit@pinyin@head:nN
              \l__hit_index_zh_tl \l__hit_index_pinyin_tl
          }
        \tl_if_empty:NF \l__hit_index_pinyin_tl
          {
            \use:x
              {
                \exp_not:N \sindex [ china ]
                  {
                    \exp_not:V \l__hit_index_pinyin_tl @
                    \exp_not:V \l__hit_index_zh_tl
                  }
              }
          }
      }
    \tl_if_empty:NF \l__hit_index_en_tl
      {
        \use:x
          { \exp_not:N \sindex [ english ] { \exp_not:V \l__hit_index_en_tl } }
      }
  }
\cs_new_protected:Npn \__hit_index_read:n #1
  {
    \tl_clear:N \l__hit_index_zh_tl
    \tl_clear:N \l__hit_index_en_tl
    \tl_clear:N \l__hit_index_pinyin_tl
    \tl_if_in:nnTF {#1} { = }
      { \hit@keys@set:nn { hit / index / one } {#1} }
      {
        \__hit_abbr_cjk_set:n {#1}
        \bool_if:NTF \l__hit_abbr_cjk_bool
          { \tl_set:Nn \l__hit_index_zh_tl {#1} }
          { \tl_set:Nn \l__hit_index_en_tl {#1} }
      }
  }
\NewDocumentCommand \hitindex { s m }
  {
    \__hit_index_read:n {#2}
    \__hit_index_emit:
    \IfBooleanF {#1}
      {
        \bool_if:NTF \g__hit_en_bool
          {
            \tl_if_empty:NTF \l__hit_index_en_tl
              { \tl_use:N \l__hit_index_zh_tl }
              { \tl_use:N \l__hit_index_en_tl }
          }
          {
            \tl_if_empty:NTF \l__hit_index_zh_tl
              { \tl_use:N \l__hit_index_en_tl }
              { \tl_use:N \l__hit_index_zh_tl }
          }
      }
  }
%    \end{macrocode}
% \end{macro}
%
% \cs{abbr} 用一次就登记一次索引，索引形式直接取条目的 \option{zh} 与
% \option{en}，不另设字段——开了 \option{indexing} 两边本来就该是同一个词。
%    \begin{macrocode}
\cs_new_protected:Npn \__hit_abbr_use:nn #1#2
  {
    \bool_set_eq:NN \l__hit_index_on_bool \g__hit_abbr_indexing_bool
    \hit@keys@set:nn { hit / abbr / use } {#1}
    \bool_if:NT \l__hit_index_on_bool
      {
        \tl_clear:N \l__hit_index_pinyin_tl
        \tl_set:Nx \l__hit_index_zh_tl { \hitabbrzh {#2} }
        \tl_set:Nx \l__hit_index_en_tl { \hitabbren {#2} }
        \__hit_index_emit:
      }
  }
%    \end{macrocode}
%
% \begin{macro}{\listofabbreviations}
% 排缩略语表。带星号是不排章标题，接在符号表那一页下面用，
% 与 \cs{listoffigures} 一族同名同姓。
%    \begin{macrocode}
\NewDocumentCommand \listofabbreviations { s O{} }
  {
    \__hit_symbols_opt:n {#2}
    \IfBooleanF {#1}
      {
        \hit@frontlist@open:nn
          { \hit@term@list@of@abbreviations@heading@zh } { \hit@term@list@of@abbreviations@heading@en }
      }
    \IfBooleanT {#1}
      {
        \tl_if_empty:NF \l__hit_symbols_subheading_tl
          { \hit@frontlist@subheading:V \l__hit_symbols_subheading_tl }
      }
%    \end{macrocode}
%
% 选项要先展开再交给 \pkg{glossaries}：\cs{glsxtrabbrvtype} 是个宏，原样传过去
% 那边拿它当类型名去查表，查不到就静悄悄排一张空表，不报错也不警告。
%
% 章标题由这里自己排（要进目录、要走部件的闸门），所以把 \pkg{glossaries} 那侧的
% 标题压掉。只给 \option{title} 一个空值不够——它照样会调 \cs{glossarysection}，
% 那是个 \cs{chapter*}，顺手把页清了，表就掉到下一页去。
%    \begin{macrocode}
    \group_begin:
    \RenewDocumentCommand \glossarysection { o m } { }
    \use:x
      {
        \exp_not:N \printnoidxglossary
          [
            type = \glsxtrabbrvtype , sort = word ,
            style = \bool_if:NTF \g__hit_en_bool { engstyle } { hitabbrstyle } ,
            title = { } , nopostdot , nonumberlist
          ]
      }
    \group_end:
  }
%    \end{macrocode}
% \end{macro}
%
% \begin{macro}{\hit@structure@merged@symbols:}
% 合并成一页的那个部件：一个章标题，符号在上、缩略语在下，两段各带小标题。
%    \begin{macrocode}
\cs_new_protected:Npn \hit@structure@merged@symbols:
  {
    \hit@structure@part:nn { list-of-symbols-and-abbreviations }
      {
        \bool_gset_false:N \g__hit_frontlist_open_bool
        \tl_clear:N \l__hit_symbols_subheading_tl
        \hit@frontlist@open:nn
          { \hit@term@list@of@symbols@and@abbreviations@heading@zh }
          { \hit@term@list@of@symbols@and@abbreviations@heading@en }
%    \end{macrocode}
%
% 章标题排过了，把闸门顶上，里头的 \env{symbols} 就不再排自己的那一个。
% 两段各预置一个小标题：符号那段取符号表的标题，缩略语那段显式给。
%    \begin{macrocode}
        \bool_gset_true:N \g__hit_frontlist_open_bool
        \tl_gset:Nn \g__hit_frontlist_pending_tl
          { \hit@term@list@of@symbols@heading@zh }
        \exp_args:Nv \input { g__hit_structure_list-of-symbols-and-abbreviations_tl }
        \listofabbreviations * [ subheading = \hit@term@list@of@abbreviations@heading@zh ]
        \bool_gset_false:N \g__hit_frontlist_open_bool
      }
  }
%    \end{macrocode}
% \end{macro}
%    \begin{macrocode}
%</final-glossary>
%    \end{macrocode}
%
% \Finale
